\documentclass{article}
\usepackage[letterpaper, left=0.75in, right=0.75in, top=1in, bottom=1in]{geometry}
\usepackage[dvipsnames]{xcolor}
\usepackage{multicol, fancyhdr, hyperref, parskip, enumitem}

\parskip=0.6em
\setlist{noitemsep, topsep=0.2em}

%%% \sub{Title}
\newcommand{\sub}[1]{
    {\raggedleft
        \large{\textsc{\color{OliveGreen}#1}}\par
    }
}

%%% \split{Date}{Title}
\newcommand{\split}[3]{
\noindent\begin{tabular}[t]{@{}l}
    \textbf{#1} \\ #2
\end{tabular}
\hfill
\begin{tabular}[t]{l@{}}
    \\
    \textit{#3}
\end{tabular}
\noindent\rule{\columnwidth}{0.5pt}
}

%%% Header
\pagestyle{fancy}
\fancyhf{}
\lhead{\large{\textbf{7-Session Intermediate Ceramics}}}
\rhead{\color{OliveGreen}\small{April 6 -- May 18, 2026}}
\rfoot{\small{Last updated \today.}}


\begin{document}

\sub{Description}

This ceramics course covers both technical skills for throwing and aesthetic
analysis in the area of pottery. We will cover a set of fundamental throwing
skills for weights up to 5~lbs with a focus on efficiency. We will also cover
analytical skills that will help you identify areas for improvement in your own
throwing as well as increase your ability to plan your work in advance. Lastly,
we set aside time for special techniques that students in the class wish to
pursue. Optionally, there will be regular sessions where we will learn how to
critique work using observational skills.

\sub{Calendar}

\begin{multicols}{2}

\split{Session 1}{Fundamentals: Water \& Centering}{4/6 7:30-9:30pm}
\begin{itemize}
    \item \textbf{Techniques}: Reducing water usage; Centering
    \item \textbf{Analysis}: Observation: the root of analysis, when to do it?
    \item \textbf{Exercises}: Only add water throwing; Slip only centering
\end{itemize}

\split{Session 2}{Fundamentals: Opening \& Pulling}{4/13 7:30-9:30pm}
\begin{itemize}
    \item \textbf{Techniques}: Opening tall instead of short; How to pull: hand position, tool/pot size ratio; Collaring with both hands; Rib shaping
    \item \textbf{Analysis}: Picking out elements of a curve: endpoints, inflections; Shape, color, and surface observations
    \item \textbf{Exercises}: Compress base with fist, raising walls in process; Small plates off hump to practice right hand control
\end{itemize}

\split{Critique 1}{Mugs}{4/17 4:00-6:00pm}
\begin{itemize}
    \item \textbf{Critique}: Bring ten mug shapes, no firing necessary.
      Separate into two groups of five based on what you think are the best or
      worst.
    \item \textbf{Exercise}: Analysing based on a rubric.
\end{itemize}

\split{Session 3}{Tall Throwing: Wall Work}{4/20 7:30-9:30pm}
\begin{itemize}
    \item \textbf{Techniques}: Centering to the top; Thick wall pulls; Collar pulls
    \item \textbf{Analysis}: Regulated vs.\ free workmanship; the workmanship of risk vs.\ certainty
    \item \textbf{Exercises}: Small cups off hump to practice collar pulls; General weight progression up to 3lbs
\end{itemize}

\columnbreak

\split{Session 4}{Tall Throwing: Big Pulls \& Planning}{4/27 7:30-9:30pm}
\begin{itemize}
    \item \textbf{Techniques}: Big pulls; Planning pots: 2/3 for body and 1/3 for neck; shaping in early pulls; cleaning bowl bottoms
    \item \textbf{Analysis}: Poorly made vs.\ well made work; Rough vs.\ precise work
    \item \textbf{Exercises}: General weight progression up to 4lb; Wire off buttons
\end{itemize}

\split{Critique 2}{Bottles}{5/1 4:00-6:00pm}
\begin{itemize}
    \item \textbf{Critique}: Bring ten bottle shapes, separate into two groups
      of five by best and worst.
    \item \textbf{Exercise}: Creating a rubric.
\end{itemize}

\split{Session 5}{Review: Putting It All Together}{5/4 7:30-9:30pm}
\begin{itemize}
    \item \textbf{Analysis}: Joining rough and precise work
    \item \textbf{Exercises}: General weight progression up to 5lb; Vase to plate to vase
\end{itemize}

\split{Session 6}{Open Studio}{5/11 7:30-9:30pm}
This is a free session where we can cover special topics requested by the
students in class. If no topics are brought forward we will cover glazing or
trimming off the hump.

\split{Critique 3}{Jars}{5/15 4:00-6:00pm}
\begin{itemize}
    \item \textbf{Critique}: Bring ten jar shapes, separate into two groups
      of five by best and worst.
    \item \textbf{Exercise}: Creating a rubric.
\end{itemize}

\split{Session 7}{Open Studio}{5/18 7:30-9:30pm}
This is a free session where we can cover special topics requested by the
students in the class. If no topics are brought forward we will cover teapots or
textures.

\end{multicols}
\end{document}
